% Ad Familiares 4.14 (ad-familiares-04-14)
% Source: https://raw.githubusercontent.com/PerseusDL/canonical-latinLit/master/data/phi0474/phi056/phi0474.phi056.perseus-lat1.xml
% Fetched by scripts/fetch_latin.py

% Dateline (Perseus): Scr. Romae in. a. 708 (46) .

\ciceroLetterOpener{M. CICERO S. D. CN. PLANCIO.}

\ciceroSection{1} binas a te accepi litteras Corcyrae datas; quarum alteris mihi gratulabare, quod audisses me meam pristinam dignitatem obtinere, alteris dicebas te velle, quae egissem, bene et feliciter evenire. ego autem, si dignitas est bene de re publica sentire et bonis viris probare quod sentias, obtineo dignitatem meam; sin autem in eo dignitas est, si, quod sentias, aut re efficere possis aut denique libera oratione defendere, ne vestigium quidem ullum est reliquum nobis dignitatis, agiturque praeclare, si nosmet ipsos regere possumus, ut ea, quae partim iam adsunt, partim impendent, moderate feramus; quod est difficile in eius modi bello, cuius exitus ex altera parte caedem ostentet, ex altera servitutem.

\ciceroSection{2} quo in periculo non nihil me consolatur, cum recordor haec me tum vidisse, cum secundas etiam res nostras, non modo adversas pertimescebam videbamque, quanto periculo de iure publico disceptaretur armis quibus si ii vicissent, ad quos ego pacis spe, non belli cupiditate adductus accesseram, tamen intellegebam, et iratorum hominum et cupidorum et insolentium quam crudelis esset futura victoria, sin autem victi essent, quantus interitus esset futurus civium partim amplissimorum, partim etiam optimorum, qui me haec praedicentem atque optime consulentem saluti suae malebant nimium timidum quam satis prudentem existimari.

\ciceroSection{3} quod autem mihi de eo, quod egerim, gratularis, te ita velle certo scio; sed ego tam misero tempore nihil novi consili cepissem, nisi in reditu meo nihilo meliores res domesticas quam rem publicani offendissem. quibus enim pro meis immortalibus beneficiis carissima mea salus et meae fortunae esse debebant, cum propter eorum scelus nihil mihi intra meos parietes tutum, nihil insidiis vacuum viderem, novarum me necessitudinum fidelitate contra veterum perfidiam muniendum putavi. sed de nostris rebus satis vel etiam nimium multa.

\ciceroSection{4} de tuis velim ut eo sis animo, quo debes esse, id est ut ne quid tibi praecipue timendum putes. si enim status erit aliquis civitatis, quicumque erit, te omnium periculorum video expertem fore; nam alteros tibi iam placatos' esse intellego, alteros numquam iratos fuisse. de mea autem in ite voluntate sic velim iudices, me, quibuscumque rebus opus esse intellegam, quamquam videam, qui sini hoc tempore et quid possim, opera tamen et consilio, studio quidem certe rei, famae, saluti tuae praesto futurum. tu velim, et quid agas et quid acturum te putes, facias me quam diligentissime certiorem. vale .
